\section{特征学习与核方法的差异}
\label{sec:14-Deep-Learning/10-Interpretability-Mathematics/02}

做一次对照实验。同一个卷积网络，第一组端到端训练；第二组把随机初始化的全部卷积层冻住，只训练最后那个线性头。在 CIFAR-10 上前者能到 \(93\%\) 左右，后者停在六成上下。然后把网络加宽——通道数翻倍、再翻倍——两条曲线之间的差距不但没有随着容量增加而扩大，反倒开始收窄。

这个走势不符合``参数越多越强''的直觉，却正好对上一个可以算得明白的结论：宽度趋于无穷时，梯度训练下的网络退化成一个固定核上的线性方法，中间层不再学习任何东西。于是深度模型区别于核方法的那一部分，恰恰在有限宽度里。把这条边界说清楚，才知道``网络就是核方法''这句话在什么范围内成立。

\subsection{无限宽把训练线性化成一个固定核上的回归}

在初始化 \(\theta_0\) 附近做一阶展开，
\[
f_\theta(x)\approx f_{\theta_0}(x)+\nabla_\theta f_{\theta_0}(x)^{\trans}(\theta-\theta_0),
\]
右端对 \(\theta\) 是线性的，对应的核是神经正切核
\[
K_{\theta_0}(x,x')=\nabla_\theta f_{\theta_0}(x)^{\trans}\nabla_\theta f_{\theta_0}(x').
\]
把平方损失下的梯度流写在函数值上，骨架只有三行。设训练集为 \(X\)，残差 \(r_t=f_{\theta_t}(X)-y\)，梯度流 \(\dot\theta=-\nabla_\theta L\) 给出
\[
\begin{aligned}
\dot f_{\theta_t}(X)&=\nabla_\theta f_{\theta_t}(X)^{\trans}\dot\theta\\
&=-\nabla_\theta f_{\theta_t}(X)^{\trans}\nabla_\theta f_{\theta_t}(X)\,r_t
=-K_t\,r_t,
\end{aligned}
\]
其中 \(K_t\) 是 \(n\times n\) 的核 Gram 矩阵。若 \(K_t\equiv K\) 不随时间变化，它便是一个线性常微分方程，解为 \(r_t=e^{-Kt}r_0\)。把 \(K\) 谱分解，残差在特征值 \(\lambda_i\) 对应的方向上以 \(e^{-\lambda_i t}\) 衰减——大特征值方向先被拟合，小特征值方向留到最后，这正是常说的谱偏置。

条件对账要列全，因为每一条在实践中都会被违反。第一，需要宽度趋于无穷并采用配套的参数缩放（输出层按 \(1/\sqrt{m}\) 缩放的所谓 NTK 参数化），才能保证训练期间参数相对初始化的移动量趋于零、\(K_t\) 趋于常量。第二，需要学习率足够小、时间连续，离散步长与自适应优化器都会引入偏离。第三，平方损失使方程线性，换成交叉熵只能得到近似。第四，\(K\) 依赖初始化的随机抽取，无穷宽时它收敛到一个确定的极限核，有限宽时它本身是随机的。这些条件成立时，网络的行为等价于以 \(K\) 为核的核回归——训练找到的是在该再生核希尔伯特空间中范数最小的插值解，几何完全由初始化定死。核与内积几何的基础见\hyperref[sec:13-Machine-Learning/09-Kernel-Methods-and-Gaussian-Processes/02]{``核技巧：只通过内积看世界''}与\hyperref[sec:13-Machine-Learning/09-Kernel-Methods-and-Gaussian-Processes/03]{``再生核希尔伯特空间：把函数当成点''}。

\subsection{NTK 解释了一批现象，也解释掉了它自己}

这个极限带来的收获是实在的。它给出深度网络为什么``容易训练''的一个回答：无穷宽下的目标在函数值上是凸的，梯度流收敛到全局插值解，不存在坏的局部极小，加宽确实让优化变简单。它给出学习曲线的形状：低频、光滑的成分先被拟合，高频成分后被拟合，这与实测的训练曲线吻合得不错。它还提供了一个可计算的基线——把 NTK 算出来做核回归，得到的数字是``完全不学表示能达到什么水平''的参照。

收获的另一面同样明确。在这个极限里，\(\nabla_\theta f_{\theta_t}\) 不随训练改变，意味着中间层给出的特征映射从头到尾就是初始化时那一个。表示没有被学过，只是被随机抽了一次然后固定下来。因此，凡是依赖``表示随任务而变''的现象，NTK 原则上都解释不了：预训练模型迁移到新任务、稀疏方向被逐步选出、深层特征比浅层更贴合语义、同一个主干支撑多个下游头——这些恰好是深度学习在实践中被使用的主要理由。

实测的差距也支持这个划分。在标准视觉基准上，NTK 核回归与同架构的端到端训练之间通常有几个到十几个百分点的差距，网络越接近实用配置差距越明显。加宽会缩小差距，因为加宽正是往极限靠；但没人是靠无限加宽把模型做好的。所以更合适的读法是，NTK 描述的是一个特征学习被关掉的对照组，它的价值在于把``网络的好表现里有多少不需要学表示''这个量估出来。

\subsection{固定几何与可训练几何的分工不同}

核方法与深度网络都可以拆成``先选表示，再在表示上做预测''。差别不在于线性还是非线性——通用核诱导的函数类同样很大，用表达能力区分两者说不通，这一点与通用逼近定理的处境相似，见\hyperref[sec:14-Deep-Learning/01-Representation-and-Approximation/02]{``通用逼近定理说清了什么又没说什么''}。差别在于几何是不是训练的自变量。

\begin{center}
\begin{tikzpicture}[>=stealth]
\begin{scope}
\foreach \p in {(0.4,1.5),(1.1,0.6),(0.5,-0.6),(1.6,-1.3),(2.3,0.2)} {
  \fill \p circle (2.4pt);
}
\foreach \p in {(2.6,1.5),(1.9,1.1),(0.2,0.3),(2.9,-0.9),(1.3,-1.6)} {
  \draw \p circle (2.6pt);
}
\draw[very thick] (0.1,-1.9) -- (3.1,1.9);
\node[below,font=\small,align=center] at (1.6,-2.2) {固定核：\\ 训练只转动这条分界线};
\end{scope}
\begin{scope}[shift={(6.2,0)}]
\foreach \p in {(0.35,1.0),(0.8,0.3),(0.25,-0.4),(0.95,-1.1),(0.6,0.9)} {
  \fill \p circle (2.4pt);
}
\foreach \p in {(2.5,1.0),(2.85,0.2),(2.4,-0.5),(3.0,-1.2),(2.7,0.7)} {
  \draw \p circle (2.6pt);
}
\draw[very thick] (1.7,-1.9) -- (1.7,1.9);
\draw[->,gray] (1.35,1.5) -- (0.85,1.2);
\draw[->,gray] (2.0,-1.6) -- (2.45,-1.35);
\node[below,font=\small,align=center] at (1.7,-2.2) {可训练表示：\\ 训练同时挪动这些点};
\end{scope}
\end{tikzpicture}
\end{center}

\emph{同一批样本，一边只能调整分界，一边可以改写样本之间的距离。}

图里左右两侧的点是同一批数据。固定核把每个样本送到 \(\phi(x)\) 处就不再移动，学习发生在这个几何里选一条分界；可训练表示允许 \(\phi_\theta\) 随损失一起更新，样本之间的距离、可分方向和局部不变性都跟着变。卷积事先写入了局部性与参数共享，注意力事先写入了内容相关的读取，训练再从这些被允许的结构里挑一个适合当前目标的表示——预先规定一部分几何，把剩下的交给数据。

这也解释了为什么``核方法表示不了、网络表示得了''不是实践差异的来源。拉开距离的因素另有几项：哪一类函数在当前表示下复杂度低；样本量增长时计算与存储怎么涨（核方法要处理 \(n\times n\) 的 Gram 矩阵，代价见\hyperref[sec:13-Machine-Learning/09-Kernel-Methods-and-Gaussian-Processes/04]{``核岭回归与核方法的代价''}）；结构先验与数据形态是否匹配；学到的相似性在分布变化之后还剩多少意义。比较两类方法时，参数量、训练预算、数据增强和核的选择都要对齐，否则单次准确率的差异说明不了发生过特征学习。

\subsection{特征有没有在动，是可以测的}

既然分界线画在``表示动不动''上，就应该把它测出来而不是猜。最直接的量是核漂移：在训练前后各算一次经验 NTK，比较两个 Gram 矩阵的对齐程度；对齐度接近 1 说明当前训练制度靠近线性化区域，明显下降说明特征在移动。第二个量是线性探测：冻结第 \(l\) 层，在其输出上训练一个线性分类器，比较训练前后的准确率随层数的变化曲线。第三个是对照组：把完整训练、只训最后一层、以及 NTK 核回归放在同一个数据与预算下比，三条曲线之间的间距就是特征学习贡献的量。第四个是宽度扫描：宽度增加时若三条曲线互相靠拢，说明你正在把模型推向那个不学表示的极限。

这些诊断都不是因果证明，报告时要留住这一点。线性探测变好可能来自信息被重新组织，也可能只是维度或正则化变了；核对齐度下降可能来自表示移动，也可能来自尺度漂移。要得到可用的结论，需要写清对照组是什么、干预的是哪个环节。泛化数字与解释之间的距离，见\hyperref[sec:13-Machine-Learning/10-Learning-Theory/04]{``泛化界的解释力与边界''}。

\subsection{什么时候该把几何固定下来}

样本少、领域知识已经给出可靠的相似性度量、需要凸优化的可复现性或者需要干净的不确定性表达时，固定核往往更省数据、更稳定、也更容易审计——它的假设全写在核里，看得见。数据与算力充足、结构先验能写进架构、并且有足够的迁移评测撑住结论时，端到端学表示的收益才兑现得出来。

所以要问的不是选传统方法还是深度方法，而是相似性里有多少应当由人预先规定，又有多少证据支持把它交给当前这批数据去学。而一旦交出去，就得回答学到的表示长什么样。有一个流传很广的答案是：数据集中在高维空间中的一片低维结构上，网络学的就是这片结构的坐标。下一节检查这个说法能被怎样验证，以及验证不了的时候该怎么讲话。
